% FY27 ONR YIP informal brief topic description. Compile using pdflatex.
% Scientific framing tailored for William Krebs (ONR Command and Control). Program Officer line can be updated after consultation.
\documentclass[11pt,letterpaper]{article}
\usepackage[letterpaper,top=0.60in,bottom=0.58in,left=0.72in,right=0.72in]{geometry}
\usepackage{mathptmx}
\usepackage[T1]{fontenc}
\usepackage{microtype}
\usepackage{enumitem}
\usepackage{xcolor}
\usepackage[colorlinks=true,citecolor=blue,linkcolor=black,urlcolor=blue]{hyperref}
\pagestyle{empty}
\setlength{\parindent}{0pt}
\setlength{\parskip}{3pt}
\setlist[itemize]{leftmargin=1.13em,topsep=2pt,itemsep=2pt,parsep=0pt}
\emergencystretch=1.5em
\newcommand{\head}[1]{\textbf{#1}\enspace}
\newcommand{\POline}{Dr. William Krebs, Code 311 --- Command and Control}
\newcommand{\FocusArea}{Command, Control, Communications, Computers, Cyber, Intelligence, Surveillance, and Reconnaissance (C5ISR) and Naval Space --- Intelligence, Decision Superiority, Command and Control (C2), and Combat Systems}
\begin{document}
\begin{center}
{\fontsize{14}{16}\selectfont\bfseries Convergent Event Reasoning for Distributed Command and Control\par}
\vspace{2pt}
{\small Brief topic description for Program Officer discussion --- ONR FY27 YIP (N0001426SF004)}
\end{center}
\vspace{-2pt}
{\small \textbf{Xinya Du}, Assistant Professor of Computer Science, The University of Texas at Dallas\par
\textbf{Email:} \href{mailto:xinya.du@utdallas.edu}{xinya.du@utdallas.edu} \quad \textbf{Website:} \href{https://xinyadu.github.io}{https://xinyadu.github.io}\par
\textbf{Program Officer:} \POline\par
\textbf{Proposed ONR S\&T Focus Area:} \FocusArea\par}
\vspace{2pt}\hrule\vspace{5pt}
\head{Problem.} Naval command nodes reason from reports that arrive late, out of order, or not at all. An afloat node may first receive a warning about an unidentified vessel and later learn it is scheduled support traffic; a shore node may receive those reports in reverse order. Even after receiving the same messages, their AI assistants may produce conflicting event histories because earlier context shapes how each report is interpreted. Data synchronization can merge records, but cannot by itself reconcile uncertain identities and event relations extracted from language.
\par\head{Central hypothesis.} We hypothesize that separating source evidence from revisable event hypotheses enables distributed agents to reconcile knowledge despite different message arrival orders. Order-independent exchange preserves the same underlying reports, while explicit uncertainty about identity, timing, and corrections supports consistent reinterpretation; prioritizing evidence that changes command questions makes limited communication more useful.
\par\head{Research plan.}
\begin{itemize}
\item \textbf{Build source-linked, revisable event knowledge.} Convert each report into source-stamped assertions about actors, events, times, locations, and relations. Record corrections as links between assertions rather than overwrites, and preserve competing identity matches when two descriptions may refer to the same event. Separate stable report-level evidence from hypotheses inferred using surrounding context, so later information can change an interpretation without changing what was originally reported.
\item \textbf{Reconcile interpretations across nodes.} Develop order-independent, duplicate-safe merging of source assertions, followed by structured reconciliation of uncertain entity links, event relations, and corrections. Align alternative interpretations using shared source references and temporal constraints rather than accepting each node's first extraction as final. Characterize when nodes holding the same evidence can recover compatible event hypotheses, and when ambiguity requires several possibilities to remain open.
\item \textbf{Share the evidence that changes command decisions.} Represent commanders' critical information requirements as standing queries over event hypotheses. Estimate which unsent assertions or corrections would change a recipient's answer, resolve a material contradiction, or reduce uncertainty, accounting for the recipient's partial view. Prioritize these updates under limited bandwidth, and generate answers with source links and explicit unresolved questions rather than transmitting full summaries indiscriminately.
\end{itemize}
\head{Evaluation and expected outcomes.} Partition \href{https://aclanthology.org/2022.emnlp-main.60/}{\textcolor{blue}{MAVEN-ERE}}, KAIROS-style event scenarios, and \href{https://proceedings.neurips.cc/paper_files/paper/2024/hash/e560a0b22e4432003d0dba63ff8dc457-Abstract-Datasets_and_Benchmarks_Track.html}{\textcolor{blue}{MEQA}} questions across nodes with delayed, duplicated, and reordered reports. Examine recovery of compatible event interpretations, accuracy of standing-query answers, and communication cost against local summarization, chronological overwrite, and full synchronization.
\par\head{Naval relevance.} During distributed maritime operations, afloat and shore nodes may receive conflicting identifications and status updates for the same vessel. Reconciling those reports and sharing evidence that changes critical information requirements could support a coherent common tactical picture across intermittent links, directly addressing uncertainty and out-of-sequence information in naval command and control.
\par\head{Prior work and PI qualifications.} Our QA-based event extraction (\href{https://doi.org/10.18653/v1/2020.emnlp-main.49}{\textcolor{blue}{Du and Cardie, 2020}}), MEQA benchmark (\href{https://openreview.net/forum?id=snNuvAOQxB}{\textcolor{blue}{Li et al., 2024}}), and DARPA KAIROS \href{https://doi.org/10.18653/v1/2022.naacl-demo.7}{\textcolor{blue}{RESIN-11}} system work provide foundations for extracting and reasoning over event knowledge. The PI holds a 2024 NSF CAREER award on consistent event graphs.

\end{document}
